\section{Introduction}\label{sec:introduction}

Let $(R,\m)$ be a commutative Noetherian regular local ring of dimension
$d$. For finite $R$-modules $M,N$ whose tensor product has finite length,
Serre's intersection multiplicity is the integer
\[
  \chi^R(M,N)=\sum_{i=0}^{d}(-1)^i
       \length_R\Tor_i^R(M,N).
\]
The support condition makes every term finite, and regularity makes the
higher Tor groups vanish. This Euler characteristic measures the local
intersection of the supports while accounting for higher Tor. Its
positivity is not a formal consequence of the nonzero degree-zero term:
the higher terms enter with alternating signs.

\begin{theorem}\label{thm:main}
Let $(R,\m)$ be a commutative Noetherian regular local ring of dimension
$d$, and let $M,N$ be nonzero finitely generated $R$-modules. If
$\length_R(M\otimes_R N)<\infty$ and
$\dim_R M+\dim_R N=d$, then
\[
  \chi^R(M,N)>0.
\]
\end{theorem}

Here regularity means that $\m$ can be generated by $d$ elements, and
module dimension means the Krull dimension of its support. Theorem
\ref{thm:main} proves Serre's positivity conjecture without a restriction
on characteristic or ramification. In particular, it does not require
either module to lift to an unramified regular local ring.\footnote{A
public post by Editan, dated December~9, 2025 and credited
to Editan, Copilot, and Gemini, also asserts the prime-quotient form of
Serre's positivity conjecture. We record this as a same-scope claim,
not as an established theorem; see
\href{https://note.com/editan/n/n35dec2c1ddbe}{the author's post}.}

\subsection{Background and significance}

Serre introduced intersection multiplicities through the Euler
characteristic of Tor. His reduction to the diagonal proves positivity
in equicharacteristic and in the unramified setting; the same argument
also treats formal power series rings over an arbitrary complete
discrete valuation ring \cite[Chapter~V]{Serre}. Thus ramification alone
does not separate the known cases from the general problem.
Roberts and, independently, Gillet--Soul\'e proved vanishing when the
tensor product has finite length and the sum of dimensions is smaller
than the ring dimension \cite{Roberts,RobertsRennes,GilletSoule}.
Their arguments use local Chern characters and Adams operations,
respectively. Gabber proved nonnegativity in the remaining
mixed-characteristic case using alterations and an intersection
calculation on a normal bundle
\cite[Section~6.1, Theorem~6.1.1]{Gabber}. Vanishing removes
lower-dimensional contributions in the reductions below. For
complementary-dimensional pairs, strict positivity must additionally
rule out zero.

Further progress in ramified mixed characteristic includes Skalit's
positivity theorem for regular local rings essentially smooth over a
two-dimensional regular local base \cite{Skalit}. KC and Soto Levins
prove positivity for modules satisfying a Serre-liftability condition
\cite[Theorem~A]{KCSotoLevins}. In their positivity statement,
$R=S/(f)$ for an unramified regular local ring $S$ and
$f\in\m_S\setminus\m_S^2$, and $M$ has a Serre lift $L$:
a finite $S$-module with $M=L\otimes_S R$ and
$\dim S-\dim L=\dim R-\dim M$; it need not satisfy the usual
Tor-vanishing condition for a lift. Theorem~\ref{thm:main} imposes
neither this liftability condition nor the smoothness hypothesis.

The proof combines normalized-length methods with algebraic $K$-theory.
Faltings introduced the normalized-length formalism in almost
mathematics \cite[Appendix~2]{Faltings}; Gabber--Ramero developed its
almost-ring-theoretic foundations further \cite{GabberRamero}.
The perfectoid length results of Cai--Lee--Ma--Schwede--Tucker
\cite{CLMST}, together with their characteristic-$p$ counterparts
\cite{Ma}, give algebras on which parameter Koszul complexes are exact
in positive degrees after passing to normalized length. The mixed-characteristic
construction uses perfectoidization and almost purity, as developed by
Bhatt--Scholze \cite{BhattScholze}. Land--Tamme's
descent theorem for truncating invariants \cite{LT} supplies the
categorical descent input. Bhatt--Hochster--Ma show that the existence
of lim Cohen--Macaulay sequences for both quotient domains implies
positivity \cite[Theorem~6.1]{BHM}. Ishiro--Shimomoto obtain such
sequences of algebras from perfectoid towers over complete
mixed-characteristic Noetherian local domains with $F$-finite residue field
\cite[Theorem~3.17]{IshiroShimomoto}. These are conditional routes
through finite modules. Here no such tower on either quotient domain
is assumed: the proof uses bounded complexes modulo free sequences
of negligible normalized rank.

\subsection{Proof strategy}

After standard reductions, it is enough to consider two complete local
domains $D=R/P$ and $E=R/Q$ of positive dimensions $h,l$, where
$h+l=d$ and $P+Q$ is $\m$-primary. Each domain has a finite extension,
denoted $D'$ or $E'$, with a finite Galois action over a regular
normalization base $A_D$ or $A_E$, that is, a regular local subring
over which the extension is finite. Associated algebras $C_D,C_E$ carry additive
nonnegative lengths $\lambda_D,\lambda_E$. Their parameter quotients
have positive length, whereas positive parameter Koszul homology has
length zero.

There are two distinct comparisons in the argument. First, rational
$K$-theory with support at the closed point shows that ordinary and
normalized Euler lengths agree after summing over all Galois
conjugates. Write $G_D,G_E$ for the two groups, and
$a_D=\rank_D D'$, $a_E=\rank_E E'$. For $\sigma\in G_D$, let
$c_\sigma$ be the normalized Tor Euler characteristic of $C_D$ and
$E'$ over $R$, where the $R$-action on $C_D$ is twisted by $\sigma$.
The first comparison identifies the sum of the $c_\sigma$ with
$|G_D|a_Da_E\chi^R(D,E)$. It therefore suffices to show that
each $c_\sigma$ is positive. Fix one such twist.

To analyze that value, choose a regular sequence $x$ in $Q$ that is a
parameter system on $D$. The module $U=C_D/xC_D$ has positive finite
normalized length over $B=R/(x)$. It can be approximated by finite-length
$B$-modules $M_n$ killed by one fixed power of the maximal ideal of $B$,
with ordinary lengths
divided by integers $s_n$ converging to the relevant normalized lengths.
Their minimal resolutions have ranks $o(s_n)$ above degree $l$.
Truncating at degree $l$ therefore produces supported complexes over
$E'$ in the category that kills free sequences whose ranks divided by
$s_n$ tend to zero.

The second comparison takes place in this category. Its supported
objects have nonnegative homology lengths, even when support appears
only after passing to the quotient. The two Euler evaluations again
agree on sums of conjugates. Evaluation using $C_E$ leaves only degree
zero, whose length is bounded below by a positive constant times the
normalized number of generators of $M_n$. The common annihilator and
the positive limiting length of $M_n$ make this lower bound strictly
positive. The second Galois average proves positivity for each value
in the first average, completing the argument.

Figure~\ref{fig:two-comparisons} records the two interfaces. Here
$b_{\sigma,\gamma}$ denotes the $C_E$-evaluation of the truncated
sequence after the additional twist $\gamma\in G_E$. Choose $a$ with
$\m_B^aM_n=0$ for all $n$, and write $\beta_0(M_n)$ for the minimal
number of generators of $M_n$. Then
$\length_B M_n\leq\beta_0(M_n)\length_B(B/\m_B^a)$.
Thus the positive normalized length of the $M_n$ forces a positive
normalized number of generators; it is this uniform estimate that
makes every $b_{\sigma,\gamma}$ strictly positive.

\begin{figure}[htbp]
\centering
\renewcommand{\arraystretch}{1.5}
\begin{tabular}{c}
\fbox{\strut Ordinary supported comparison over $A_D\subset D'$}\\
$\displaystyle |G_D|a_Da_E\,\chi^R(D,E)
       =\sum_{\sigma\in G_D}c_\sigma$\\
$\Big\uparrow$\quad prove positivity separately for each $\sigma$\\
\fbox{\strut Asymptotic supported comparison over $A_E\subset E'$}\\
$\displaystyle |G_E|c_\sigma
       =\sum_{\gamma\in G_E}b_{\sigma,\gamma}>0$
\end{tabular}
\caption{Two norm comparisons, not invariance of individual evaluations.
The first is Lemma~\ref{pos:first-average}; the second is
Equation~\eqref{pos:second-average}. The uniform annihilator and
generator bound give the strict inequality in
Lemma~\ref{pos:positive-evaluation}.}
\label{fig:two-comparisons}
\end{figure}

The reusable technical points are the finite-surjection descent
argument with arbitrary connective algebra coefficients and the
construction of supported Euler evaluations after quotienting by
negligible ranks. In particular, the comparison is proved on
idempotent summands in the quotient, not only on complexes that already
have uniform support before taking the quotient. No invariance of a
normalized-length function under either Galois group is assumed.

Section~\ref{sec:lengths} establishes the precise normalized-length
inputs. Section~\ref{sec:descent} proves supported rational descent and
the first Euler comparison. Section~\ref{sec:asymptotic} constructs the
rank quotient and its supported evaluations. Section~\ref{sec:positivity}
performs the approximation and the two comparisons to prove
Theorem~\ref{thm:main}.

\subsection{Consequences}

Theorem~\ref{thm:main} also has consequences for symbolic powers.
For a ramified mixed-characteristic regular local ring $(R,\m)$ and
primes $P,Q$ with $\sqrt{P+Q}=\m$ and
$\operatorname{ht}P+\operatorname{ht}Q=\dim R$, the implication of
Kurano--Roberts \cite[Theorem~3.2]{KuranoRoberts} now gives
\[
 P\cap Q^{(r)}\subseteq\m^{r+1}\qquad(r\geq1),
 \qquad Q^{(r)}:=Q^rR_Q\cap R.
\]
Here ramified means that the residue characteristic $p$ lies in
$\m^2$. This is a direct application of their positivity-to-symbolic-power
theorem, rather than an additional ingredient in our proof.

The next consequence applies Skalit's finite-support blowup argument
\cite[Corollary~D]{Skalit} with Theorem~\ref{thm:main} as its positivity
input. Dutta studied the lower bound in this finite-support setting
\cite{DuttaFiniteSupport}. We include the argument to exhibit how
strict positivity at the points of the blowup controls equality.

\begin{corollary}[Finite strict-transform intersection]\label{cor:tangent-cone}
Let $(A,\m)$ be a commutative Noetherian regular local ring with residue
field $k$, and let $P,Q$ be prime ideals such that $A/(P+Q)$ has finite
length and $\dim(A/P)+\dim(A/Q)=\dim A$. Write $e_\m$ for Hilbert--Samuel
multiplicity and put
\[
 T=\operatorname{gr}(A/P)\otimes_{\operatorname{gr}A}\operatorname{gr}(A/Q),
\]
where all associated graded rings use the maximal-ideal filtrations.
If $\dim T\leq1$, then
\[
 \chi^A(A/P,A/Q)\geq e_\m(A/P)e_\m(A/Q),
\]
with equality if and only if $\dim T=0$.
\end{corollary}

\begin{proof}
If $\dim A=0$, then $A=k$, $P=Q=0$, and both sides are $1$ with $T=k$.
Otherwise, let $\widetilde X$ be the blowup of $X=\Spec A$ at its closed
point, and let $\widetilde Y,\widetilde Z$ be the strict transforms of
$Y=\Spec(A/P)$ and $Z=\Spec(A/Q)$. By \cite[Lemma~3.5]{SkalitUnramified},
the underlying set $S$ of $\widetilde Y\cap\widetilde Z$ identifies with
that of $\operatorname{Proj}T$. Thus $S$ is a finite set of closed points
of the exceptional divisor, empty exactly when $\dim T=0$.
Fulton's blowup formula and Skalit's finite-support calculation
\cite[Theorem~3.4 and proof of Proposition~3.10]{SkalitUnramified} give
\[
 \chi^A(A/P,A/Q)-e_\m(A/P)e_\m(A/Q)
   =\sum_{x\in S}[\kappa(x):k]\,\chi^{B_x}(M_x,N_x),
\]
where $B_x=\mathcal O_{\widetilde X,x}$,
$M_x=\mathcal O_{\widetilde Y,x}$, and $N_x=\mathcal O_{\widetilde Z,x}$.
Each residue degree is finite and positive since the exceptional divisor
is projective over $k$. The blowup $\widetilde X$ is regular
\cite[Proposition~A.4(1)]{Jannsen}. The modules $M_x,N_x$ are nonzero and
finite over $B_x$, with finite-length tensor product because their
intersection is supported on $S$. The local domains $A/P$ and $A/Q$,
as well as $A$, are formally equidimensional and universally catenary.
Thus \cite[Lemma~3.8 and proof of Proposition~3.10]{SkalitUnramified}
gives $\dim M_x=\dim(A/P)$, $\dim N_x=\dim(A/Q)$, and
$\dim B_x=\dim A$. Theorem~\ref{thm:main} therefore makes every local
intersection multiplicity in the sum strictly positive. The asserted
inequality follows, and equality holds exactly when $S$ is empty.
\end{proof}

\subsection{Conventions}

We fix universes large enough for the module and sequence categories
under consideration.
Complexes are indexed homologically. An associative derived algebra is
\emph{connective} when its homotopy groups in negative degrees vanish.
Perfect modules over associative algebras are right modules.
All scalar localizations of coefficient algebras are along the central
image of the commutative base. We use nonconnective algebraic
$K$-theory of perfect complexes; a subscript $\Q$ denotes rationalization.
For a sequence $a=(a_1,\ldots,a_r)$, write $K(a;V)$ for its homological
Koszul complex with coefficients in $V$, and put
$H_i(a;V)=H_i(K(a;V))$. The notation $a^v$ denotes
$(a_1^v,\ldots,a_r^v)$, not the $v$th power of the ideal it generates.
Normalized lengths are allowed to vanish on nonzero modules; all uses
of their additivity involve finite values or nonnegative extended
values for which the indicated operations are defined.
